Notably, our transformer-based classifier achieves a remarkable accuracy of 94.3\% on the held-out test set, significantly outperforming the CNN baseline (89.7\%) \citep{george2018, zevin2017}. This represents a substantial leap forward in glitch classification. The model was trained on $2.1\times10^5$ labeled spectrograms from the Gravity Spy dataset \citep{zevin2017}, leveraging a 12-layer architecture with 8 attention heads — a design choice that proved pivotal. Interestingly, performance on the ``Blip'' and ``Koi Fish'' classes was notably lower (F1 = 0.71 and 0.68, respectively), which may be attributable to their morphological similarity, limited training examples, and inherent label noise. It is worth noting that the attention maps delve into physically meaningful time-frequency regions, underscoring the interpretability of the approach. In summary, these results highlight the transformative potential of attention-based architectures for gravitational-wave data quality — a crucial step toward real-time detector characterization \citep{cabero2019}.
